% 맺음
% ====================================================================
\section*{Part T 맺음}
\addcontentsline{toc}{section}{Part T 맺음}

본 파트는 코인 하프셀의 엔트로피와 가역 발열을 \emph{통계열역학}으로 세웠다. 분배함수 $Z$ 에서 점유 분포
$\avg{n}$ 을 끌어내 그것이 Part I logistic(\S\ref{sec:width})의 기원임을 보였고($w=RT/F$ 가 분포 폭), 점유 분포에서
configurational 엔트로피와 그 발산하는 부분몰 항을, 같은 분포 언어로 vibrational(포논 Bose--Einstein)$\cdot$
electronic(Fermi--Dirac--Sommerfeld) 엔트로피를 세웠으며, 진동 항의 온도 편차를 Einstein 모드 보정으로 따로
닫았다. 세 분포가 합쳐 측정 $\partial U_\oc/\partial T(x)$ 의 모든 모양을 만든다: 겹침 가중(파생 A, 수치 검증
완료)으로 봉우리 사이 연속 블렌드를, 봉우리 내부 config(파생 B, 이중계산 분리)로 $x$-의존을, 폭 $w_j$ 의
이중지위(파생 C --- 단상은 평형 예측, 흑연 두-상은 현상학적 자유 피팅 폭, broadening 기원은 Part I
\S\ref{sec:broadening})로
두-상 폭의 지위를, 히스 분기 평균(파생 D)으로 가역/비가역 분리를 닫았다. 극한 여섯 코너가 분포 식을 검산하고,
``전이당 상수 $\Delta S^0_j$ $+$ 분포''로 측정급 비선형이 모델 안에서 자기일관하게 산출됨을 보였다(지배 두-상
전이의 폭 $T$-의존은 다온도 round-trip 으로 확정할 대상). 한 계산 예제($\bar x{=}0.25$, \S\ref{ssec:worked})가
이 종합식~\eqref{eq:complete} 를 구체 수치로 닫아 유한차분 round-trip 과 $<0.001\,\mu$V/K 로 일치시켰고, 다섯
SOC 로 확장한 표~\ref{tab:qrev} 가 저-$\bar x$ 발열$\cdot$고-$\bar x$ 흡열의 부호 교대를 같은 식 하나로
재현했다 --- 다온도 $dQ/dV$ 에서 $\Delta S^0_j$ 를 얻는 절차(방법론 출구 박스)도 함께 제시했다. 가역 발열
$\dot Q_\rev=-IT\,\partial U_\oc/\partial T$ 은, 서두에서 목적함수로 노출한 그대로, 이 세 분포를 재배열하는
열이다.

여기까지가 Part T 흑연 열특성의 결과 사슬이다 --- 곡선(Part I)과 열특성(Part T)의 전 입력 인자를 한자리에
모으는 진행 요약이 §18(\S\ref{sec:inputs})로 이어지고, 부호 사슬 검산표$\cdot$코드 구현 요구명세 부록들이
그 뒤를 잇는다.

% 자산: [C-117] 맺음 4파생 역할(A 연속블렌드·B config 이중계산분리·C 폭 이중지위·D 히스 분기평균 가역/비가역) ·
%   [C-118] 6코너 분포식 검산·"상수+분포"로 측정급 자기일관(두-상 폭 T-의존 round-trip 확정) ·
%   [C-119] 가역발열=분포 재배열 열(서두 C-6과 수미상관)
